\chapter{Verification commands}
\label{app:verify}

The checks that discriminate between a working policy and a dead one. Ordered by
how often they were the one that mattered.

\section{Is the policy actually live?}

\begin{lstlisting}[style=term]
condor_status -af Name Start OwnerSessionActive HostMemAvailMB HostDiskAvailMB
\end{lstlisting}

Reads the \textbf{slot ad}. If a policy attribute is missing here, the policy is
dead no matter what the configuration says. Allow one update interval, or force
a fresh ad with \kn{condor\_reconfig}.

\section{Why is this machine not taking work?}

\begin{lstlisting}[style=term]
condor_status -af Name Start TotalLoadAvg HostCpuPsiAvg10 \
                       HostMemAvailMB HostDiskAvailMB
\end{lstlisting}

Every term in \kn{START} beside its threshold. A machine can show free slots and
still decline work --- that is the design, not a fault.

\section{What did the startd actually decide?}

\begin{lstlisting}[style=term]
grep -iE 'suspend|continue|changing activity|changing state' \
     /var/log/condor/StartLog | tail -25
\end{lstlisting}

The only place claim refusals are recorded. \kn{Request to claim resource
refused} and \kn{Claiming protocol failed} appear here and nowhere else.

\section{Is the pool intact?}

\begin{lstlisting}[style=term]
condor_status -af Name State Activity TotalLoadAvg HostMemAvailMB Start
condor_q -totals
\end{lstlisting}

\section{Did the configuration take?}

\begin{lstlisting}[style=term]
condor_config_val -v <KNOB>
\end{lstlisting}

\begin{trap}
Necessary but far from sufficient. Every defect in Chapter~\ref{ch:failure}
printed a correct value here. The \kn{-v} flag at least reveals \emph{which} file
in \file{config.d} won, which is its real use --- precedence bites.
\end{trap}

\section{Observing under load}

\begin{lstlisting}[style=term]
cut -d' ' -f1 /proc/loadavg          # the real figure, not the collector's
awk '/^MemAvailable:/ {print int($2/1024)" MB"}' /proc/meminfo
cat /proc/pressure/cpu /proc/pressure/memory
\end{lstlisting}

\begin{trap}
Do not use \kn{condor\_status} to observe a stress test: it reads the collector
and is only as fresh as \kn{UPDATE\_INTERVAL}. It showed a frozen value for two
and a half minutes while load was driven from 1 to 15.

\kn{condor\_status -direct} is not an alternative on a worker --- it needs the
pool signing key to verify the caller's token, and only the entry node has it.
\end{trap}

\section{Load generation}

Both generators must be self-terminating, with \kn{timeout} owning the deadline
so nothing survives a dropped session.

\begin{lstlisting}[style=term]
%% CPU: tight loop, no fork per iteration
timeout "$SECS" sh -c 'while :; do :; done'

%% Memory: touch one byte per page so the pages are genuinely resident
timeout "$SECS" python3 -c '
import sys, time
mb = int(sys.argv[1])
buf = bytearray(mb * 1024 * 1024)
for i in range(0, len(buf), 4096):
    buf[i] = 1
time.sleep(int(sys.argv[2]))
' "$MB" "$SECS"
\end{lstlisting}

An untouched allocation is virtual only, and \kn{MemAvailable} would not move ---
the test would appear to fail while never having applied any pressure.

\section{Testing eviction without endangering the machine}

Raise the floor to meet the memory rather than driving the machine into genuine
pressure. Override \textbf{both} expressions:

\begin{lstlisting}
PREEMPT      = (HostMemAvailMB < <just under current availability>)
WANT_SUSPEND = !(HostMemAvailMB < <the same value>)
\end{lstlisting}

Overriding \kn{PREEMPT} alone leaves \kn{WANT\_SUSPEND} on the real floor, where
it stays true, so \kn{SUSPEND} is consulted instead and the test fails for the
wrong reason.
